% Digital Report Template
% Uses doc-digital for dark, screen-optimized premium output
% Compile with: lualatex main.tex

\documentclass[11pt]{article}

% Core digital style (includes fonts, colors, dark bg, layout, TikZ)
\usepackage{doc-digital}

% Optional: math notation shortcuts
% \usepackage{doc-math}

% Optional: additional table environments
% \usepackage{doc-tables}

% Document metadata
\title{Quarterly Platform Report}
\author{Pinpoint Engineering}
\date{Q1 2026}

\begin{document}

% ---- Title Page ----
\digitaltitle[Platform Health and Roadmap Summary]{Quarterly Platform Report}{Pinpoint Engineering}{March 2026}

% ---- Table of Contents ----
{
    \hypersetup{linkcolor=PinTextMuted}
    \tableofcontents
}
\newpage

% ---- Content ----
\section{Executive Summary}

This quarter saw significant progress across the platform. Test execution throughput increased by 34\% following the migration to parallel dispatchers, while the average bug report resolution time decreased from 48 hours to under 18.

\begin{keyinsight}
The combination of API token expansion and MCP server integration has reduced onboarding time for new enterprise clients by roughly half compared to the previous quarter.
\end{keyinsight}

\section{Platform Metrics}

\subsection{Reliability}

System uptime remained at 99.97\% across all production regions. The single incident on February 12 was traced to a certificate rotation failure in the staging proxy, which did not affect customer-facing services.

\begin{table}[ht]
    \centering
    \caption{Availability by Region}
    \begin{tabular}{lcc}
        \toprule
        \textcolor{PinTextBright}{\textbf{Region}} &
        \textcolor{PinTextBright}{\textbf{Uptime}} &
        \textcolor{PinTextBright}{\textbf{Incidents}} \\
        \midrule
        US East     & 99.99\% & 0 \\
        EU West     & 99.97\% & 1 \\
        APAC        & 99.98\% & 0 \\
        \bottomrule
    \end{tabular}
\end{table}

\subsection{Performance}

\begin{criticalpoint}
Peak API response time exceeded the 200ms SLA target on three occasions during batch import windows. The dispatcher queue redesign in v2.4 addresses this with priority-based routing.
\end{criticalpoint}

\section{Architecture Updates}

The platform architecture shifted toward a more modular dispatch layer this quarter. The key changes are summarized below.

\begin{itemize}
    \item Migrated from monolithic dispatcher to per-project worker pools
    \item Introduced API token scoping for CLI and MCP integrations
    \item Added tenant-aware caching at the service boundary
\end{itemize}

\begin{technicaldetail}
Worker pools use a weighted fair-queuing algorithm that allocates capacity proportional to each tenant's contracted tier. Enterprise tenants receive dedicated pool reservations, while Professional and Core tenants share a common pool with burst capacity.
\end{technicaldetail}

\subsection{Diagram: Dispatch Flow}

\begin{center}
\begin{tikzpicture}[node distance=2cm and 2.5cm]
    \node[component] (api) {API Gateway};
    \node[component, right=of api] (dispatch) {Dispatcher};
    \node[database, right=of dispatch] (db) {PostgreSQL};
    \node[api, below=of dispatch] (worker) {Worker Pool};

    \draw[dataflow] (api) -- (dispatch);
    \draw[dataflow] (dispatch) -- (db);
    \draw[dataflow] (dispatch) -- (worker);
    \draw[dataflow, dashed, PinSuccess!50] (worker) -- (db) node[midway, right, font=\sffamily\tiny, text=PinTextMuted] {results};
\end{tikzpicture}
\end{center}

\section{Looking Ahead}

Q2 priorities focus on three areas:

\begin{enumerate}
    \item \glow{Onboarding automation} through project templates and guided setup flows
    \item \glow{Reporting engine} for scheduled bug report digests per project
    \item \glow{CLI feature parity} with the web dashboard for power users
\end{enumerate}

\begin{businessvalue}
Automated onboarding is projected to reduce the average time-to-first-test from 3 days to under 4 hours, directly improving trial-to-paid conversion rates.
\end{businessvalue}

\begin{pullquote}
The best testing platform is the one your team actually uses every day.
\end{pullquote}

\begin{warningbox}
The legacy webhook format will be deprecated in v2.5. All integrations must migrate to the v2 event schema before June 30, 2026.
\end{warningbox}

\end{document}
